%----------------------------------------------------------------------
% Abkuerzungs- und Glossareintraege
%----------------------------------------------------------------------
\makenoidxglossaries
%----------------------------------------------------------------------
% --- Abkuerzungen (Abkuerzungsverzeichnis) ---
%----------------------------------------------------------------------
\newacronym{ki}{KI}{Künstliche Intelligenz}
\newacronym{llm}{LLM}{Large Language Model, großes Sprachmodell}
\newacronym{vlm}{VLM}{Vision-Language Model, Sprach-Bild-Modell}
\newacronym{mcp}{MCP}{Model Context Protocol}
\newacronym{som}{SoM}{Set-of-Mark}
\newacronym{vr}{VR}{Virtual Reality, virtuelle Realität}
\newacronym{mr}{MR}{Mixed Reality, gemischte Realität}
\newacronym{vpn}{VPN}{Virtual Private Network, virtuelles privates Netz}
\newacronym{json}{JSON}{JavaScript Object Notation}
\newacronym{api}{API}{Application Programming Interface, Programmierschnittstelle}
\newacronym{gpu}{GPU}{Graphics Processing Unit, Grafikprozessor}
\newacronym{vram}{VRAM}{Video Random Access Memory, Grafikspeicher}
\newacronym{stt}{STT}{Speech-to-Text, Spracherkennung}
\newacronym{rq}{RQ}{Research Question, Forschungsfrage}
\newacronym{fa}{FA}{funktionale Anforderung}
\newacronym{nfa}{NFA}{nicht-funktionale Anforderung}
\newacronym{es}{ES}{Execution Success, Ausführungserfolg}
%----------------------------------------------------------------------
% --- Glossarbegriffe ---
%----------------------------------------------------------------------
\newglossaryentry{lokaler-betrieb}{
  name={lokaler Betrieb},
  description={Selbst gehosteter Betrieb der Modelle ohne Cloud-Dienst,
  siehe Abschnitt~\ref{sec:ein-motivation}}
}
\newglossaryentry{strukturelle-repraesentation}{
  name={strukturelle Repräsentation},
  description={Bereitstellung der Szene als exakte Laufzeitdaten der Engine
  (Objekt-ids, Typen, Positionen, Größen, Farben, Kamerapose) in textueller
  Form, hier als JSON}
}
\newglossaryentry{visuelle-repraesentation}{
  name={visuelle Repräsentation},
  description={Bereitstellung der Szene über eine von einem Sprach-Bild-Modell
  erzeugte natürlichsprachliche Beschreibung der gerenderten, mit
  Set-of-Mark-Markern versehenen Ansicht, siehe
  Abschnitt~\ref{sec:grundlagen-modelle}}
}
\newglossaryentry{naive-hybride-repraesentation}{
  name={naive hybride Repräsentation},
  description={Kombination aus struktureller und visueller Schicht ohne
  verbindliche Instruktion zur Zuständigkeit der beiden Quellen}
}
\newglossaryentry{priorisierende-hybride-repraesentation}{
  name={priorisierende hybride Repräsentation},
  description={Hybride Repräsentation mit vorangestellter visueller
  Beschreibung, die als maßgeblich für das Aussehen deklariert wird}
}
\newglossaryentry{kompetenzgetrennte-hybride-repraesentation}{
  name={kompetenzgetrennte hybride Repräsentation},
  description={Hybride Repräsentation mit expliziter Aufgabenteilung: die
  visuelle Beschreibung bestimmt, welches Objekt gemeint ist, die
  strukturelle Schicht liefert alle numerischen Werte}
}
\newglossaryentry{kennungssystem}{
  name={gemeinsames Kennungssystem},
  description={Einheitliche Objektkennungen, über die alle
  Repräsentationsvarianten Objekte ansprechen; die visuelle Variante ist über
  ihre Bildmarken an dieselben Kennungen gebunden}
}
\newglossaryentry{grounding}{
  name={Grounding},
  description={Auflösung einer sprachlichen Referenz auf ein konkretes
  Szenenobjekt}
}
\newglossaryentry{folgezustand}{
  name={Folgezustand},
  description={Tatsächlicher Zustand der Szene nach einer Manipulation, aus
  dem der Ausführungserfolg unabhängig von der Selbstauskunft des Agenten
  bestimmt wird}
}
\newglossaryentry{referenztyp}{
  name={Referenztyp},
  description={Art, in der ein natürlichsprachlicher Befehl auf ein Objekt
  verweist, etwa über eine Kennung, ein Attribut, eine räumliche Relation,
  eine rein sichtbare Eigenschaft oder ein verdecktes Objekt}
}
